\subsection*{All-amplitude vacuum coefficient: independence audit}

\paragraph*{Recommended status.}
\textbf{Derived by 2 distinct complete methods at continuum
two-spatial-derivative bulk order}
[X: resummed polar heat current; arbitrary-background Sylvester response
with the isospectral trace identity].
Use this one status wherever the result is used.
P16 itself contains both methods.  The final3 author and RED reconstruct
them; their separate readership supplies no additional method.
Foundation item (d) displays the heat method and cites the longer
reconstructions.  Its displayed argument alone is not a second method.
% SOURCES: P16 response.md, Secs. 1 and 3; attachment
% P16_nonlinear_polar_vacuum.md, Secs. 2.1--2.5 and 4.1--4.2;
% P16_COPY_BACK.md, Main result.  Exact roots are registered below.
% Audited local author, Secs. 2--3 and Earned check tags; local RED,
% Duty 4 and Independent disposition; foundational_arguments_final3.tex (d).

\paragraph*{Shared hypotheses and different decisive steps.}
Choose the one-sided polar operator
$A_s=m_s/\psi+i\psi^{-3}D_a$ for the one full sine operator,
allocation $L=N\psi$, and squared-mass rectangle
$x_s=(x_0,x_0+h_1,x_0+h_2,x_0+h_1+h_2)$ with
$\eta=(+,-,-,+)$.
Require $x_0,h_1,h_2>0$, a smooth real lapse and a smooth commuting ruler
$0<P_-\leq\psi\leq P_+<\infty$, a compatible closed Dirac domain,
and a grouped spin trace.  The density
$X(x)=\sum_s\eta_s\operatorname{tr}_2
 (|A_s^\dagger|-|A_s|)_{xx}$ includes all sine valleys and is
independent of $N$.  In one continuum valley write
$A_s=M+f\sigma\cdot\nabla$, $M=m_s/\psi$, $f=\psi^{-3}$.

The heat route computes the two partner connections and potentials,
then uses scalar and Pauli insertions in the resummed diagonal heat
kernel.  Their difference is a divergence; its grouped square-root
integral gives the coefficient.  The author and RED supply the
radial-gauge argument excluding an order-two curvature insertion.
That supplies a derivation of the heat step, but remains the same
method as P16's heat calculation.

The second route uses no such heat coefficient.  Its local expansion
can be organized directly with, for either positive polar square $H$,
\[
 \sqrt H=\frac2\pi\int_0^\infty H(H+t^2)^{-1}\,dt .
\]
Expanding the resolvent symbol in slow spatial derivatives produces
local functions of the ruler jets: its leading continuum inverse is
$(M^2+f^2p^2+t^2)^{-1}$, and higher symbol-product terms carry
further spatial derivatives.  The positive gap controls the
infrared, and the grouped mass moments control the coefficient's
ultraviolet limit.  This is the gapped derivative expansion used in
both proofs, not an assumption of the desired coefficient or a
uniform remainder bound on an arbitrary profile.
% SOURCES: local author, Smooth domain and an operator gap, and
% Why arbitrary uniform P fixes the two-derivative function; P16
% attachment Secs. 1.3--1.4 and 4.2.  The scalar resolvent identity
% follows by integrating lambda/(lambda+t^2) for lambda>0; its
% symbol expansion explains locality without the heat-current identity.

\paragraph*{Closure of the second route.}
Cubic symmetry allows only
$X^{[2]}=u(\psi)\Delta\psi+v(\psi)|\nabla\psi|^2$.
The exact finite-regulator singular-value identity is
$\operatorname{Tr}(|A^\dagger|-|A|)=0$.
Isolating its two-derivative term on arbitrary slow periodic profiles
gives
\[
 0=\int [v(\psi)-u'(\psi)]|\nabla\psi|^2 .
\]
For $\psi=P+\delta\phi$, the $\delta^2$ coefficient forces
$v(P)=u'(P)$ for every $P>0$.  Thus one coefficient function suffices;
the trace identity alone does not determine it.

At each uniform $P$, set $\delta\psi=\phi$.  Pauli multiplication and differentiation
of $E_j^2=P_j$, where $P_0=A^\dagger A$, $P_1=AA^\dagger$, give
\[
 \begin{aligned}
 \operatorname{tr}_2(\delta P_1-\delta P_0)_{p+q,p}
   &=-6P^{-7}|d_a(p+q)-d_a(p)|^2\phi_q,\\
 (\delta E_j)_{p+q,p}
   &=\frac{(\delta P_j)_{p+q,p}}{e_{p+q}+e_p},\\
 u_a(P)&=\frac3{P^4}\sum_s\eta_s\int_{\rm BZ}
 \frac{d^3p}{(2\pi)^3}
 \frac{\cos^2(ap_1)}{\sqrt{|d_a(p)|^2+x_sP^4}},
 \end{aligned}
\]
where $d_{a,i}=a^{-1}\sin(ap_i)$ and
$e_p=P^{-3}\sqrt{|d_a(p)|^2+x_sP^4}$ for each packet.
This is the response at every positive background, not just $P=1$.

The continuum integral can be evaluated without a heat representation.
The rectangle is the double mass derivative, so in one valley
\[
 \begin{aligned}
 \mathcal I(P)&:=\int\frac{d^3p}{(2\pi)^3}
       \sum_s\eta_s(p^2+x_sP^4)^{-1/2}\\
 &=\frac{3P^8}{4}\int_0^{h_1}\!du\int_0^{h_2}\!dv
   \int\frac{d^3p}{(2\pi)^3}
       [p^2+P^4(x_0+u+v)]^{-5/2}\\
 &=\frac{P^4J}{8\pi^2},\qquad
 J=\int_0^{h_1}\!du\int_0^{h_2}\!dv\,
       (x_0+u+v)^{-1}>0 .
 \end{aligned}
\]
The radial integral is $1/(6\pi^2c)$ for denominator
$(p^2+c)^{5/2}$.  Counting the 8 valleys of the one sine operator
once gives $u(P)=3J/\pi^2=:k$, hence $v=u'=0$.  Therefore
\[
 X^{[2]}=k\Delta\psi,\qquad
 W_{0,\mathrm{bulk}}^{[2]}
   =-k\int\nabla\psi\cdot\nabla(N\psi).
\]
The coefficient, its absence of additional amplitude functions, and
its normalization follow without the heat-current identity.
This closes a different complete method for this claim.
% SOURCES: P16 attachment Secs. 4.1--4.2, Eqs. (28)--(31);
% local author Sec. 3; local RED Duty 4, A separate arbitrary-background
% coefficient check.  The positive double-difference radial evaluation
% above makes explicit the heat-free evaluation of the shared integral.

\paragraph*{Limits of the status.}
The shared operator, grading, regulator, gap and local expansion are
hypotheses, not independent physical evidence.  The result is exact
in amplitude at two derivatives only.  At finite spacing $u_a(P)$
runs with $aP^2$; a single reference normalization does not make it
constant.  On a region with boundary retain
$k\oint L\partial_n\psi$ and the domain's boundary terms.
Neither method supplies the complete nonlocal functional, a
massless uniform remainder, or a stellar approximation bound.
No numerical check or [N] tag is used for this independence judgment;
[R] is not assigned to this reviewer B derivation.

\paragraph*{Exact source locators.}
\begingroup\footnotesize
Protected P16 root:
\path{<internal>/}.
Files: \path{response.md}, sections 1 and 3;
\path{files/P16_nonlinear_polar_vacuum.md}, sections 1.3--2.5
and 4.1--4.3, and section 9 for the nonlocal scope;
\path{files/P16_COPY_BACK.md}, \emph{Main result}.
Local audit root:
\path{<internal>/}.
Files: \path{rederivation_final3_vacuum_surface.tex}, sections 2--3
and \emph{Earned check tags and exact limits};
\path{redteam_final3_vacuum_surface.tex}, duty 4 and
\emph{Independent disposition and exact check tags};
\path{foundational_arguments_final3.tex}, item (d).
\endgroup
% Governing SOURCES read:
% <internal>/GOAL_alt.md, duty 2;
% <internal>/GOAL_lane_z.md;
% <internal>/JUDGMENT_web40c.md,
% Secs. 12--13, and NOTES_shift_and_spacing_window.md;
% <internal>/DECISIONS_ADDENDUM.md, (g)--(r);
% <internal>/JUDGMENT_gaps2.md;
% <internal>/STATE.tex and DECISIONS.md.
